\chapter{Contribuții}

Înainte de construirea rețelei, a fost nevoie de un set de date care să conțină informații prelucrate extrase din segmente
 audio ce conțin doar elemente vocale, izolate de partea instrumentală sau fișiere audio ce permit izolarea părții
 instrumentale de cea vocală și apoi prelucrarea părții vocale. Seturi de date deja existente nu au putut fi utilizate
 deoarece nu conțin fișierele audio, ci doar date prelucrate din care nu se pot izola elementele vocale sau, în cazul
 celor care conțin fișiere audio, nu au un număr suficient de date sau o diversitate destul de mare a genurilor muzicale.
 Așadar, a fost necesară construirea unui set de date etichetate format din fișiere audio ce conțin doar elementele vocale
 din melodiile originale.

A fost construit atat un pipeline de procesare a setului de date și salvarea lui sub forma unor reprezentări ușor de
 preluat ca input pentru rețeaua neuronală, cat și un pipeline de încărcare în memorie a acestor reprezentări. În ceea
 ce privește rețeaua neuronală, au fost făcute experimente cu mai multe tipuri de arhitecturi în urma cărora a fost
 selectată arhitectura cu cele mai bune rezultate, o rețea ce primește ca input spectograme Mel, utilizează o rețea
 convoluțională pentru a reduce dimensionalitatea spectrogramei și a extrage caracteristicile ei esentiale, aceste
 caracteristici reprezentând inputul pentru o rețea recurentă care tratează aspectul temporalității spectrogramei.

În final, a fost construită aplicația web care permite vizualizarea predicțiilor pe date noi, în comparație cu segmentele
 din setul de date cu vectori de predicție similari. 
